\section{Mathematical Framework}

\subsection{Statistical models, parameters, and estimators}

\begin{definition}[Model]
A statistical model is a family of probability measures
\[
\mathcal P=\{P_\theta:\theta\in\Theta\}
\]
on a measurable sample space $(\mathcal X,\mathcal A)$.
Unless stated otherwise, $X_1,\ldots,X_n$ are independent and
identically distributed with common law $P_{\theta_0}$, where
$\theta_0$ is the true parameter.
\end{definition}

\begin{definition}[Identifiable Parametrization]
The parametrization is \emph{identifiable} if
\[
P_\theta=P_{\theta'}
\quad\Longrightarrow\quad
\theta=\theta'.
\]
Identifiability concerns the entire distribution. It does not imply
that a particular collection of moments identifies the parameter.
\end{definition}

\begin{definition}[Estimator and Statistical Functional]
An estimator of a target $q(\theta)\in\mathbb R^k$ is a measurable
function
\[
\widehat q_n:\mathcal X^n\longrightarrow\mathbb R^k.
\]
Before observing the data, $\widehat q_n$ is a random variable.
After observing $x_1,\ldots,x_n$, its value is a numerical estimate.

A statistical functional is a map
\[
\nu:\mathcal P\longrightarrow\mathbb R^k.
\]
A target $q(\theta)$ defines such a functional through
$\nu(P_\theta)=q(\theta)$ whenever this value is well defined:
\[
P_\theta=P_{\theta'}
\quad\Longrightarrow\quad
q(\theta)=q(\theta').
\]
\end{definition}
\begin{definition}[Empirical Measure]
The empirical measure is
\[
P_n=\frac1n\sum_{i=1}^n\delta_{X_i}.
\]
For an integrable function $g$, write
\[
Pg=\int g\,dP,
\qquad
P_ng=\frac1n\sum_{i=1}^ng(X_i).
\]
\end{definition}

\subsection{Basic probabilistic notation}

For deterministic $a_n>0$,
\[
Z_n=o_P(a_n)
\quad\Longleftrightarrow\quad
\frac{\|Z_n\|}{a_n}\xrightarrow{P}0.
\]
Also,
\[
Z_n=O_P(a_n)
\]
means that $Z_n/a_n$ is bounded in probability: for every
$\varepsilon>0$, there exist $M<\infty$ and $N$ such that
\[
P(\|Z_n\|>Ma_n)<\varepsilon,
\qquad n\geq N.
\]
In particular, convergence in distribution to a finite random
vector implies $O_P(1)$.

The main tools used below are:
\begin{enumerate}
\item The law of large numbers:
if $E\|Y_1\|<\infty$, then
\[
\overline Y_n\xrightarrow{P}EY_1.
\]

\item The multivariate central limit theorem:
if $E\|Y_1\|^2<\infty$, then
\[
\sqrt n(\overline Y_n-\mu)
\xrightarrow{d}N_d(0,\Sigma),
\qquad
\mu=EY_1,\quad \Sigma=\operatorname{Cov}(Y_1).
\]

\item The continuous mapping theorem:
if $Z_n\xrightarrow{d}Z$ and $g$ is continuous at $Z$ almost surely,
then
\[
g(Z_n)\xrightarrow{d}g(Z).
\]

\item Slutsky's theorem:
if $Z_n\xrightarrow{d}Z$ and $U_n\xrightarrow{P}c$, where $c$ is
constant, then
\[
(Z_n,U_n)\xrightarrow{d}(Z,c).
\]
Continuous transformations, sums, products, and ratios with a
nonzero limiting denominator may then be applied.
\end{enumerate}

\section{Construction of Estimators}

\subsection{Minimum contrast estimation}

Let $\rho:\mathcal X\times\Theta\to\mathbb R$ be a measurable
contrast function. Define the population and empirical objectives by
\[
M_{\theta_0}(\theta)=E_{\theta_0}\rho(X_1,\theta),
\qquad
M_n(\theta)=P_n\rho(\cdot,\theta).
\]
The population objective should identify the truth:
\[
M_{\theta_0}(\theta)>M_{\theta_0}(\theta_0),
\qquad \theta\neq\theta_0.
\]
A minimum contrast estimator is a measurable choice
\[
\widehat\theta_n\in
\operatorname*{arg\,min}_{\theta\in\Theta}M_n(\theta).
\]

This definition involves separate questions:
\begin{enumerate}
\item Does a minimizer exist?
\item Is it unique?
\item Can a minimizer be selected measurably?
\item Does it converge to $\theta_0$?
\end{enumerate}
None of these follows merely from writing an $\operatorname{arg\,min}$.

A continuous objective attains its minimum on a compact parameter
space. On a noncompact space, a useful sufficient condition is
compactness of suitable sublevel sets. Strict convexity on a convex
parameter space gives uniqueness whenever a minimizer exists.

If $\Theta$ is open and $M_n$ is differentiable, every interior
minimizer satisfies
\[
\nabla_\theta M_n(\widehat\theta_n)=0.
\]
The converse requires additional conditions. A stationary point may
be a local maximum or a saddle point, and boundary minimizers need
not satisfy this equation.

\subsection{Estimating equations}

An estimating function is a measurable map
\[
\psi:\mathcal X\times\Theta\longrightarrow\mathbb R^d.
\]
An estimating-equation estimator is a measurable solution of
\[
P_n\psi(\cdot,\widehat\theta_n)=0.
\]
The corresponding population equation is
\[
E_{\theta_0}\psi(X_1,\theta)=0.
\]
For identification by this equation, its unique solution should be
$\theta_0$.

Minimum contrast estimation is a special case when
\[
\psi(x,\theta)=\nabla_\theta\rho(x,\theta).
\]
Conversely, an estimating equation need not arise from a scalar
contrast function.

Terminology varies: estimators defined by minimizing empirical
objectives are commonly called \emph{M-estimators}, while roots of
empirical estimating equations are often called \emph{Z-estimators}.
The asymptotic theory of a smooth interior M-estimator usually
proceeds through its estimating equation.

\subsection{Method of moments}

Let $\theta\in\Theta\subseteq\mathbb R^d$, and suppose
$g_1,\ldots,g_d$ are measurable functions with finite expectations.
Define the moment map
\[
m(\theta)=
\begin{pmatrix}
E_\theta g_1(X_1)\\
\vdots\\
E_\theta g_d(X_1)
\end{pmatrix},
\qquad
\widehat m_n=
\begin{pmatrix}
P_ng_1\\
\vdots\\
P_ng_d
\end{pmatrix}.
\]
The method of moments solves
\[
m(\widehat\theta_n)=\widehat m_n.
\]
If $m$ is injective and $\widehat m_n\in m(\Theta)$, then
\[
\widehat\theta_n=m^{-1}(\widehat m_n).
\]

The classical choice is $g_j(x)=x^j$, giving
\[
m_j(\theta)=E_\theta X_1^j,
\qquad
\widehat m_{n,j}=\frac1n\sum_{i=1}^nX_i^j.
\]
These are raw moments, not centered moments.

Several conditions must be distinguished:
\begin{enumerate}
\item The chosen expectations must exist.
\item The chosen moment map must identify the parameter.
\item The observed moment vector must lie in the range of that map.
\item The inverse, or selected solution, must be measurable.
\item Continuity of the inverse is needed for the usual consistency
argument.
\end{enumerate}

If the target is only $q(\theta)$ and
\[
q(\theta)=h(m(\theta)),
\]
one may estimate it by
\[
\widehat q_n=h(\widehat m_n)
\]
without recovering the entire parameter.

\paragraph{Example: Bernoulli model.}
For $X_i\sim\operatorname{Bernoulli}(\theta)$,
\[
P_\theta(X_i=1)=\theta,
\qquad
P_\theta(X_i=0)=1-\theta.
\]
Since $E_\theta X_i=\theta$,
\[
\widehat\theta_n=\overline X_n.
\]

\paragraph{Example: Gamma model.}
Use the shape--rate parametrization
\[
f_{\alpha,\lambda}(x)
=
\frac{\lambda^\alpha}{\Gamma(\alpha)}
x^{\alpha-1}e^{-\lambda x},
\qquad x>0,\quad \alpha,\lambda>0.
\]
Then
\[
EX=\frac{\alpha}{\lambda},
\qquad
\operatorname{Var}(X)=\frac{\alpha}{\lambda^2}.
\]
Writing
\[
\widehat v_n=P_nX^2-\overline X_n^2,
\]
moment matching gives, when $\widehat v_n>0$,
\[
\widehat\alpha_n=\frac{\overline X_n^2}{\widehat v_n},
\qquad
\widehat\lambda_n=\frac{\overline X_n}{\widehat v_n}.
\]
The denominator is $n$, rather than $n-1$, because empirical
moments are being matched directly.

\subsection{Plug-in and extension principles}

If a target is a functional $\nu(P)$ defined at the empirical
measure, its empirical plug-in estimator is
\[
\widehat\nu_n=\nu(P_n).
\]
Examples include
\[
\nu(P)=PX,
\qquad
\nu(P_n)=\overline X_n,
\]
and
\[
\nu(P)=PX^2-(PX)^2,
\qquad
\nu(P_n)=P_nX^2-\overline X_n^2.
\]

For a distribution function $F$, define the lower quantile by
\[
F^{-1}(u)=\inf\{x:F(x)\geq u\},
\qquad 0<u<1.
\]
An empirical quantile is obtained by replacing $F$ with $F_n$.
Quantile conventions can differ at flat portions or jumps; the
chosen convention should be stated.

A parameter may initially be defined only on a model
$\mathcal P_0$, while $P_n\notin\mathcal P_0$. To apply the plug-in
principle, extend it to a larger class:
\[
\overline\nu:\mathcal P_1\to\mathbb R^k,
\qquad
\overline\nu|_{\mathcal P_0}=\nu.
\]
Then use $\overline\nu(P_n)$.

Such extensions need not be unique. Two functionals can agree
throughout the model and disagree away from it, producing different
finite-sample estimators.

For a finite sample space $\{x_1,\ldots,x_k\}$, write
\[
p_j=P(X=x_j),
\qquad
\widehat p_j=\frac1n\sum_{i=1}^n1_{\{X_i=x_j\}}.
\]
A frequency plug-in estimator of $h(p_1,\ldots,p_k)$ is
\[
h(\widehat p_1,\ldots,\widehat p_k).
\]

\subsection{Least squares}

Consider fixed design points $z_1,\ldots,z_n$ and observations
\[
Y_i=g(\beta_0,z_i)+\varepsilon_i,
\qquad E\varepsilon_i=0.
\]
Least squares minimizes
\[
Q_n(\beta)=\sum_{i=1}^n
\{Y_i-g(\beta,z_i)\}^2.
\]
If $\operatorname{Var}(\varepsilon_i)=\sigma_i^2<\infty$, then
\[
EQ_n(\beta)
=
\sum_{i=1}^n\sigma_i^2
+
\sum_{i=1}^n
\{g(\beta_0,z_i)-g(\beta,z_i)\}^2.
\]
Thus $\beta_0$ minimizes the population objective. It is the unique
minimizer precisely when the design-specific mean map
\[
\beta\longmapsto
\bigl(g(\beta,z_1),\ldots,g(\beta,z_n)\bigr)
\]
is injective.

Normal errors are not required for this population argument.

If $g$ is differentiable, an interior least-squares solution satisfies
\[
\sum_{i=1}^n
\nabla_\beta g(\widehat\beta,z_i)
\{Y_i-g(\widehat\beta,z_i)\}=0.
\]

\subsection{Linear least squares and generalized least squares}

In the linear model
\[
Y=Z\beta_0+\varepsilon,
\qquad
Z\in\mathbb R^{n\times d},
\]
the normal equations are
\[
Z^\top Z\widehat\beta=Z^\top Y.
\]
If $\operatorname{rank}(Z)=d$, the solution is unique:
\[
\widehat\beta_{\mathrm{OLS}}
=
(Z^\top Z)^{-1}Z^\top Y.
\]
If $Z$ is rank deficient, minimizers still exist, but the coefficient
vector is not unique. The fitted vector $Z\widehat\beta$ is unique.

Under
\[
E\varepsilon=0,
\qquad
\operatorname{Cov}(\varepsilon)=\sigma^2I_n,
\]
\[
E\widehat\beta_{\mathrm{OLS}}=\beta_0,
\qquad
\operatorname{Cov}(\widehat\beta_{\mathrm{OLS}})
=
\sigma^2(Z^\top Z)^{-1}.
\]
The Gauss--Markov theorem states that OLS has minimum covariance,
in the positive-semidefinite order, among linear unbiased
estimators.

More generally, suppose
\[
\operatorname{Cov}(\varepsilon)=\sigma^2W,
\]
where $W$ is known and positive definite. Generalized least squares
minimizes
\[
(Y-Z\beta)^\top W^{-1}(Y-Z\beta).
\]
For full column rank $Z$,
\[
\widehat\beta_{\mathrm{GLS}}
=
(Z^\top W^{-1}Z)^{-1}Z^\top W^{-1}Y,
\]
and
\[
\operatorname{Cov}(\widehat\beta_{\mathrm{GLS}})
=
\sigma^2(Z^\top W^{-1}Z)^{-1}.
\]
This is OLS applied to the whitened model
\[
W^{-1/2}Y=W^{-1/2}Z\beta_0+W^{-1/2}\varepsilon.
\]
For independent heteroscedastic errors, $W$ is diagonal, and the
weights are inversely proportional to the error variances.

\subsection{Likelihood and maximum likelihood estimation}

Suppose $P_\theta$ has density $p_\theta$ relative to a fixed
dominating measure. For observed data $x_1,\ldots,x_n$, define
\[
L_n(\theta)=\prod_{i=1}^np_\theta(x_i),
\qquad
\ell_n(\theta)=\sum_{i=1}^n\log p_\theta(x_i),
\]
with $\log0=-\infty$.

The likelihood is a function of $\theta$ with the data fixed.
It is not a probability density over $\Theta$.

An MLE is a measurable selection
\[
\widehat\theta_n\in
\operatorname*{arg\,max}_{\theta\in\Theta}L_n(\theta).
\]
The maximizing set can be empty or contain multiple points.

For an interior differentiable maximizer,
\[
\nabla_\theta\ell_n(\widehat\theta_n)=0.
\]
Solving this equation alone does not establish global maximality.

If $\eta=q(\theta)$ is a one-to-one reparametrization, then
$q(\widehat\theta_n)$ is an MLE of $\eta$. For a general map $q$,
the induced likelihood is
\[
L_n^q(\eta)=
\sup_{\theta:q(\theta)=\eta}L_n(\theta).
\]
A global MLE $\widehat\theta_n$ yields a maximizer
$q(\widehat\theta_n)$ of this profile likelihood.

If
\[
p_\theta(x)=a_\theta(T(x))b(x),
\]
with $b(x)$ independent of $\theta$, likelihood maximization depends
on the data only through $T(x)$.

\paragraph{KL interpretation.}
For probability measures $P\ll Q$, define
\[
D_{\mathrm{KL}}(P\|Q)
=
\int\log\left(\frac{dP}{dQ}\right)dP.
\]
It is nonnegative, with equality exactly when $P=Q$.
Whenever the relevant log expectations are finite,
\[
E_{\theta_0}[-\log p_\theta(X)]
-
E_{\theta_0}[-\log p_{\theta_0}(X)]
=
D_{\mathrm{KL}}(P_{\theta_0}\|P_\theta).
\]
Thus identifiability makes the true parameter the unique population
minimizer of negative log-likelihood.

\paragraph{Boundary effects.}
For Bernoulli observations, the MLE is $\overline X_n$ if
$\Theta=[0,1]$. If $\Theta=(0,1)$, no MLE exists when all observations
are zero or all are one: the likelihood supremum is approached at
an excluded boundary point.

\subsection{Exponential families}

A canonical exponential family has densities
\[
p_\eta(x)
=
h(x)\exp\{\eta^\top T(x)-A(\eta)\},
\]
where
\[
A(\eta)=
\log\int e^{\eta^\top T(x)}h(x)\,d\nu(x).
\]
The natural parameter space consists of the $\eta$ for which this
integral is positive and finite.

On its interior,
\[
\nabla A(\eta)=E_\eta T(X),
\qquad
\nabla^2A(\eta)=\operatorname{Cov}_\eta(T(X)).
\]
In a minimal representation, no nonzero linear combination of
$T(X)$ is almost surely constant. Consequently,
$\nabla^2A(\eta)$ is positive definite.

For an i.i.d.\ sample,
\[
\ell_n(\eta)
=
n\{\eta^\top\overline T_n-A(\eta)\}+C,
\qquad
\overline T_n=\frac1n\sum_{i=1}^nT(X_i).
\]
The score equation is
\[
\nabla A(\widehat\eta_n)=\overline T_n.
\]
The Hessian is
\[
\nabla^2\ell_n(\eta)=-n\nabla^2A(\eta),
\]
so the log-likelihood is strictly concave in a minimal family.
Hence any solution of the score equation is the unique MLE.

\paragraph{Existence and convex support.}
Let $C$ be the closed convex hull of the support of $T(X)$.
For a full minimal exponential family whose full natural parameter
space is open, a finite MLE exists exactly when
\[
\overline T_n\in\operatorname{int}(C).
\]
For a nonminimal representation, the appropriate geometric object
is the relative interior in the affine hull, and the natural
parameter itself is not uniquely identified.

The requirement that the full natural parameter space be open is
important; the geometric criterion must not be applied without
checking the family assumptions.

For a multinomial model, interior empirical frequencies mean that
every category has a positive count. Zero counts prevent a finite
canonical-parameter MLE, even though an MLE exists in the closed
probability simplex.

\subsection{Curved exponential families}

A curved exponential family restricts the natural parameter to
\[
\eta=c(\theta),
\qquad
\theta\in\Theta\subseteq\mathbb R^d,
\quad d<k.
\]
Its log-likelihood is
\[
\ell_n(\theta)
=
n\{c(\theta)^\top\overline T_n-A(c(\theta))\}+C.
\]
An interior maximizer satisfies
\[
Dc(\theta)^\top
\{\overline T_n-\nabla A(c(\theta))\}=0.
\]
Strict concavity in $\eta$ does not generally imply concavity in
$\theta$, or uniqueness of the maximizer over a curved constraint
set.

\section{Consistency}

\subsection{Pointwise and uniform consistency}

An estimator $\widehat q_n$ is consistent for $q(\theta)$ if
\[
P_\theta\bigl(\|\widehat q_n-q(\theta)\|>\varepsilon\bigr)
\longrightarrow0
\]
for every fixed $\theta$ and every $\varepsilon>0$.

It is uniformly consistent on $K\subseteq\Theta$ if
\[
\sup_{\theta\in K}
P_\theta\bigl(\|\widehat q_n-q(\theta)\|>\varepsilon\bigr)
\longrightarrow0.
\]
Uniform consistency is stronger: the parameter at which the error
is largest may change with $n$.

\paragraph{Example.}
The sample mean is consistent at every distribution with a finite
first absolute moment. If a class of distributions satisfies
\[
\sup_{P\in\mathcal P}\operatorname{Var}_P(X)\leq C<\infty,
\]
then Chebyshev's inequality gives
\[
\sup_{P\in\mathcal P}
P\bigl(|\overline X_n-EX|>\varepsilon\bigr)
\leq\frac{C}{n\varepsilon^2}.
\]
Thus the consistency is uniform on that class.

\subsection{Consistency of moment and plug-in estimators}

Suppose
\[
\widehat m_n\xrightarrow{P}m(\theta_0),
\qquad
q(\theta)=h(m(\theta)),
\]
and $h$ is continuous at $m(\theta_0)$. Then
\[
h(\widehat m_n)\xrightarrow{P}q(\theta_0).
\]
The estimator must be defined on an event whose probability tends
to one; it may be assigned an arbitrary fixed value elsewhere.

For method-of-moments estimation, this gives consistency when the
moment functions are integrable and the inverse moment map is
continuous at the true moment vector.

For a finite sample space, the empirical probability vector is
uniformly consistent over the probability simplex. Therefore every
continuous functional on that simplex has a uniformly consistent
frequency plug-in estimator, by compactness and uniform continuity.

\subsection{An argmin consistency theorem}

Let $M_n:\Theta\to\mathbb R$ be random objectives and
$M:\Theta\to\mathbb R$ deterministic. Suppose:
\begin{enumerate}
\item Uniform convergence:
\[
\sup_{\theta\in\Theta}|M_n(\theta)-M(\theta)|
\xrightarrow{P}0.
\]

\item Separation of the minimizer: for every $\varepsilon>0$,
\[
\inf_{\|\theta-\theta_0\|\geq\varepsilon}M(\theta)
>
M(\theta_0).
\]

\item Approximate minimization:
\[
M_n(\widehat\theta_n)
\leq \inf_{\theta\in\Theta}M_n(\theta)+r_n,
\qquad r_n=o_P(1).
\]
\end{enumerate}
Then
\[
\widehat\theta_n\xrightarrow{P}\theta_0.
\]

Indeed, if the separation gap is $\delta_\varepsilon>0$, then
outside a small-probability event the uniform approximation error
and $r_n$ are too small for an approximate minimizer to lie outside
the $\varepsilon$-ball.

If $\Theta$ is compact and $M$ is continuous with a unique
minimizer, separation follows automatically. On a noncompact space,
uniqueness alone does not imply separation.

This theorem supplies a framework, not an automatic proof:
uniform convergence and control of parameters escaping to infinity
must still be established.

\subsection{Consistency in exponential families}

In a full minimal exponential family with open natural parameter
space,
\[
\overline T_n\xrightarrow{P}\nabla A(\eta_0).
\]
The true mean lies in the interior of the convex support.
Locally, $\nabla A$ has a continuous inverse because
$\nabla^2A(\eta_0)$ is positive definite. Consequently,
\[
P_{\eta_0}(\text{a finite MLE exists})\longrightarrow1,
\qquad
\widehat\eta_n\xrightarrow{P}\eta_0.
\]
An arbitrary definition on the event of nonexistence does not
affect these asymptotic conclusions.

\section{Delta Methods and Distributional Approximations}

\subsection{First-order Delta method}

Suppose
\[
a_n(U_n-u)\xrightarrow{d}Z,
\qquad a_n\to\infty,
\]
and $g:\mathbb R^d\to\mathbb R^k$ is differentiable at $u$.
Then
\[
a_n\{g(U_n)-g(u)\}\xrightarrow{d}Dg(u)Z.
\]

The underlying expansion is
\[
g(U_n)-g(u)
=
Dg(u)(U_n-u)+o_P(a_n^{-1}).
\]
In particular, if
\[
\sqrt n(U_n-u)\xrightarrow{d}N_d(0,\Sigma),
\]
then
\[
\sqrt n\{g(U_n)-g(u)\}
\xrightarrow{d}
N_k\bigl(0,Dg(u)\Sigma Dg(u)^\top\bigr).
\]

\subsection{Degenerate first derivatives}

If $g'(u)=0$, the first-order Delta method yields only a degenerate
limit. Suppose, in one dimension,
\[
a_n(U_n-u)\xrightarrow{d}Z,
\qquad
g'(u)=0,
\qquad
g''(u)\neq0.
\]
A second-order Taylor expansion gives
\[
a_n^2\{g(U_n)-g(u)\}
\xrightarrow{d}\frac12g''(u)Z^2.
\]
Thus the correct rate and limiting distribution can change.

More generally, if the first nonzero derivative is of order $r$,
then, under the corresponding Taylor remainder condition,
\[
a_n^r\{g(U_n)-g(u)\}
\xrightarrow{d}
\frac{g^{(r)}(u)}{r!}Z^r.
\]

\subsection{Moment approximations require additional conditions}

Weak convergence does not imply convergence of moments.
Therefore an asymptotic normal distribution alone does not justify
approximations to the bias, variance, or mean squared error.

Let $\overline Y_n$ be the mean of i.i.d.\ scalar observations with
mean $\mu$ and variance $\sigma^2$. If
\[
E|Y_1|^3<\infty
\]
and $h$ has a bounded third derivative, Taylor expansion and moment
bounds give
\[
Eh(\overline Y_n)
=
h(\mu)+\frac{\sigma^2h''(\mu)}{2n}
+O(n^{-3/2}).
\]
Thus the leading bias is
\[
\operatorname{Bias}\{h(\overline Y_n)\}
=
\frac{\sigma^2h''(\mu)}{2n}+O(n^{-3/2}).
\]

If $E|Y_1|^4<\infty$ and $h''$ is bounded, then
\[
\operatorname{Var}\{h(\overline Y_n)\}
=
\frac{\sigma^2h'(\mu)^2}{n}
+O(n^{-3/2}).
\]
These are sufficient conditions, not minimal ones.

For i.i.d.\ vectors $Y_i$ with mean $\mu$ and covariance $\Sigma$,
appropriate third-moment and bounded third-derivative conditions
give
\[
Eh(\overline Y_n)
=
h(\mu)+
\frac1{2n}\operatorname{tr}\{D^2h(\mu)\Sigma\}
+O(n^{-3/2}).
\]
Under suitable fourth-moment and second-derivative bounds,
\[
\operatorname{Var}\{h(\overline Y_n)\}
=
\frac1n\nabla h(\mu)^\top\Sigma\nabla h(\mu)
+O(n^{-3/2}).
\]

\subsection{Studentization}

Suppose
\[
\sqrt n(\widehat q_n-q_0)\xrightarrow{d}N(0,V),
\qquad V>0,
\]
and $\widehat V_n\xrightarrow{P}V$. Then
\[
\frac{\sqrt n(\widehat q_n-q_0)}{\sqrt{\widehat V_n}}
\xrightarrow{d}N(0,1).
\]
Replacing an unknown scale by a consistent estimator is called
studentization.

For example, if $EX_1^2<\infty$ and $\sigma^2>0$, then
\[
\frac{\sqrt n(\overline X_n-\mu)}{S_n}
\xrightarrow{d}N(0,1).
\]
This asymptotic result does not require normal observations.
Exact Student $t$ distributions do require stronger assumptions.

\subsection{Variance-stabilizing transformations}

Suppose
\[
\sqrt n(\widehat\theta_n-\theta)
\xrightarrow{d}N(0,V(\theta)),
\qquad V(\theta)>0.
\]
A differentiable transformation $h$ stabilizes the first-order
variance if
\[
h'(\theta)^2V(\theta)=c
\]
for a constant $c>0$. An increasing choice satisfies
\[
h'(\theta)=\sqrt{\frac{c}{V(\theta)}}.
\]
Then
\[
\sqrt n\{h(\widehat\theta_n)-h(\theta)\}
\xrightarrow{d}N(0,c).
\]

Examples include:
\begin{enumerate}
\item Poisson mean:
$V(\lambda)=\lambda$, so $h(\lambda)=2\sqrt\lambda$ gives $c=1$.

\item Bernoulli probability:
$V(p)=p(1-p)$, so $h(p)=2\arcsin\sqrt p$ gives $c=1$ for $0<p<1$.

\item Correlation in a bivariate normal model:
$V(\rho)=(1-\rho^2)^2$, so
\[
h(\rho)=\operatorname{arctanh}(\rho)
=\frac12\log\frac{1+\rho}{1-\rho}
\]
gives $c=1$ for $|\rho|<1$.
\end{enumerate}
Variance stabilization and improvement of the shape of a
finite-sample distribution are distinct objectives.

\subsection{Edgeworth expansions}

Let
\[
Z_n=\frac{\sqrt n(\overline X_n-\mu)}{\sigma},
\]
and define standardized cumulants
\[
\gamma_1=\frac{E(X-\mu)^3}{\sigma^3},
\qquad
\gamma_2=\frac{E(X-\mu)^4}{\sigma^4}-3.
\]
Under sufficiently strong moment and regularity conditions,
for example a suitable higher-moment assumption together with
Cramer's condition on the characteristic function,
\[
\begin{aligned}
P(Z_n\leq x)
={}&\Phi(x)-\phi(x)
\biggl[
\frac{\gamma_1}{6\sqrt n}H_2(x)\\
&\qquad
+\frac1n\left\{
\frac{\gamma_2}{24}H_3(x)
+\frac{\gamma_1^2}{72}H_5(x)
\right\}
\biggr]+o(n^{-1}),
\end{aligned}
\]
uniformly in $x$, where
\[
H_2(x)=x^2-1,\qquad
H_3(x)=x^3-3x,\qquad
H_5(x)=x^5-10x^3+15x.
\]
The correction incorporates skewness and excess kurtosis.
Finite moments alone do not guarantee this displayed expansion;
lattice distributions generally require additional correction terms.

\section{Asymptotic Normality of Estimators}

\subsection{Estimating equations and asymptotic linearity}

Suppose $\widehat\theta_n$ satisfies
\[
P_n\psi(\cdot,\widehat\theta_n)=0,
\qquad
\widehat\theta_n\xrightarrow{P}\theta_*,
\]
where
\[
P\psi(\cdot,\theta_*)=0.
\]
Define
\[
A=P[D_\theta\psi(\cdot,\theta_*)],
\qquad
B=P[\psi(\cdot,\theta_*)\psi(\cdot,\theta_*)^\top].
\]
Assume:
\begin{enumerate}
\item $\theta_*$ is an interior point.
\item $B$ is finite and $A$ is nonsingular.
\item $\psi(x,\theta)$ is continuously differentiable near $\theta_*$
for almost every $x$.
\item For some $\delta>0$,
\[
P\left[
\sup_{\|\theta-\theta_*\|\leq\delta}
\|D_\theta\psi(X,\theta)\|
\right]<\infty.
\]
\end{enumerate}
These local smoothness and domination conditions give the required
local law of large numbers for the derivative.

Expanding along the segment from $\theta_*$ to
$\widehat\theta_n$ yields
\[
\widehat\theta_n-\theta_*
=
-A^{-1}P_n\psi(\cdot,\theta_*)+o_P(n^{-1/2}).
\]
Consequently,
\[
\sqrt n(\widehat\theta_n-\theta_*)
\xrightarrow{d}
N_d(0,A^{-1}BA^{-\top}).
\]
The matrix $A^{-1}BA^{-\top}$ is the sandwich covariance.

The function
\[
\operatorname{IF}(x;P)
=
-A^{-1}\psi(x,\theta_*)
\]
is the asymptotic influence function in this expansion:
\[
\widehat\theta_n-\theta_*
=
\frac1n\sum_{i=1}^n\operatorname{IF}(X_i;P)
+o_P(n^{-1/2}).
\]

The multidimensional formula is a useful extension of the
one-dimensional theory:
\[
V=
\frac{E\psi(X,\theta_*)^2}
{\{E\partial_\theta\psi(X,\theta_*)\}^2}.
\]

\subsection{Misspecified models}

The true law $P$ need not belong to the working model
$\{P_\theta\}$. A minimum contrast estimator can instead converge to
\[
\theta_*=\operatorname*{arg\,min}_\theta
E_P\rho(X,\theta),
\]
provided the minimizer is identified and consistency conditions
hold.

For maximum likelihood, the target is a KL projection when that
interpretation is well defined:
\[
\theta_*\in
\operatorname*{arg\,min}_\theta D_{\mathrm{KL}}(P\|P_\theta).
\]
The sandwich covariance remains relevant. Under misspecification,
the information identity generally fails, so it is usually
incorrect to replace the sandwich covariance by inverse Fisher
information.

\subsection{Score and Fisher information}

For one observation, let
\[
s_\theta(x)=\nabla_\theta\log p_\theta(x).
\]
Under conditions permitting differentiation under the integral,
\[
E_\theta s_\theta(X)=0.
\]
The Fisher information per observation is
\[
I(\theta)
=
E_\theta[s_\theta(X)s_\theta(X)^\top].
\]
Under the corresponding second-order regularity conditions,
\[
I(\theta)
=
-E_\theta[\nabla_\theta^2\log p_\theta(X)].
\]
For $n$ independent observations, information is $nI(\theta)$.

Common support and integrable derivative bounds are standard
sufficient ingredients. Parameter-dependent support can invalidate
these identities and the usual $\sqrt n$ theory.

\subsection{Regular maximum likelihood estimators}

Suppose a consistent interior MLE exists with probability tending
to one, and the preceding estimating-equation conditions hold for
the score. Then
\[
\widehat\theta_n-\theta_0
=
I(\theta_0)^{-1}
\frac1n\sum_{i=1}^ns_{\theta_0}(X_i)
+o_{P_{\theta_0}}(n^{-1/2}),
\]
and
\[
\sqrt n(\widehat\theta_n-\theta_0)
\xrightarrow{d}
N_d(0,I(\theta_0)^{-1}).
\]
This is a conditional theorem: consistency, interiority,
nonsingular information, and local smoothness must be verified.

For a canonical exponential family,
\[
s_\eta(X)=T(X)-\nabla A(\eta),
\qquad
I(\eta)=\nabla^2A(\eta).
\]
Thus
\[
\widehat\eta_n-\eta_0
=
I(\eta_0)^{-1}
\{\overline T_n-\nabla A(\eta_0)\}
+o_P(n^{-1/2}).
\]
This also follows from the Delta method applied to
$(\nabla A)^{-1}$.

\paragraph{Example: normal mean and variance.}
For $X_i\sim N(\mu,\sigma^2)$,
\[
\widehat\mu_n=\overline X_n,
\qquad
\widehat\sigma_n^2
=
\frac1n\sum_{i=1}^n(X_i-\overline X_n)^2.
\]
For $\sigma^2>0$,
\[
\sqrt n
\begin{pmatrix}
\widehat\mu_n-\mu\\
\widehat\sigma_n^2-\sigma^2
\end{pmatrix}
\xrightarrow{d}
N_2\left(
0,
\begin{pmatrix}
\sigma^2&0\\
0&2\sigma^4
\end{pmatrix}
\right).
\]

\section{Asymptotic Efficiency and Its Limits}

\subsection{An information bound for regular estimating equations}

Consider a scalar parameter and an estimating function satisfying
\[
E_\theta\psi(X,\theta)=0
\]
for all $\theta$ near the truth. Differentiating this identity under
the integral gives
\[
E_\theta\partial_\theta\psi(X,\theta)
=
-E_\theta[\psi(X,\theta)s_\theta(X)].
\]
Therefore its asymptotic variance is
\[
V_\psi(\theta)
=
\frac{E_\theta\psi(X,\theta)^2}
{\{E_\theta[\psi(X,\theta)s_\theta(X)]\}^2}.
\]
Cauchy--Schwarz implies
\[
V_\psi(\theta)\geq\frac1{I(\theta)}.
\]
Equality holds precisely when
\[
\psi(X,\theta)=a(\theta)s_\theta(X)
\quad P_\theta\text{-almost surely}
\]
for a nonzero scalar $a(\theta)$.

Thus regular score-based estimators achieve the information bound
within this class. This is not an unrestricted comparison with
every measurable estimator.

\subsection{Hodges's estimator and superefficiency}

Let $X_i\sim N(\theta,1)$ and define
\[
\widetilde\theta_n=
\begin{cases}
0,&|\overline X_n|\leq n^{-1/4},\\
\overline X_n,&|\overline X_n|>n^{-1/4}.
\end{cases}
\]
For every fixed $\theta\neq0$,
\[
P_\theta(\widetilde\theta_n=\overline X_n)\to1,
\]
so
\[
\sqrt n(\widetilde\theta_n-\theta)
\xrightarrow{d}N(0,1).
\]
At $\theta=0$,
\[
P_0(\widetilde\theta_n=0)\to1,
\qquad
\sqrt n\widetilde\theta_n\xrightarrow{P}0.
\]
The estimator is superefficient at zero.

This does not contradict the preceding information bound because
the estimator is not regular in the required local sense.
Indeed, along $\theta_n=cn^{-1/4}$ with $0<|c|<1$,
\[
P_{\theta_n}(\widetilde\theta_n=0)\to1,
\]
and
\[
nE_{\theta_n}(\widetilde\theta_n-\theta_n)^2
\geq
c^2\sqrt n\,
P_{\theta_n}(\widetilde\theta_n=0)
\longrightarrow\infty.
\]
Pointwise improvement can therefore conceal poor nearby behavior.

\subsection{Growing parameter dimension: an incidental-parameter example}

Suppose
\[
X_{ij}\sim N(\mu_i,\sigma^2),
\qquad
i=1,\ldots,p,\quad j=1,\ldots,k,
\]
independently, with fixed $k\geq2$ and $p\to\infty$.
The MLEs are
\[
\widehat\mu_i=\overline X_i,
\qquad
\widehat\sigma^2
=
\frac1{pk}\sum_{i=1}^p\sum_{j=1}^k
(X_{ij}-\overline X_i)^2.
\]
Since each group residual sum has mean $(k-1)\sigma^2$,
\[
\widehat\sigma^2\xrightarrow{P}
\frac{k-1}{k}\sigma^2.
\]
Thus the MLE is inconsistent, while
\[
\frac1{p(k-1)}
\sum_{i=1}^p\sum_{j=1}^k
(X_{ij}-\overline X_i)^2
\]
is consistent.

Fixed-dimensional MLE theory does not apply when the number of
nuisance parameters grows with the sample size.

\section{One-Dimensional Asymptotic Testing}

\subsection{Level, power, and consistency}

A possibly randomized test is a measurable function
\[
\varphi_n:\mathcal X^n\to[0,1].
\]
Its power at $\theta$ is $E_\theta\varphi_n$.

For testing
\[
H_0:\theta\leq\theta_0
\qquad\text{against}\qquad
H_1:\theta>\theta_0,
\]
uniform asymptotic level at most $\alpha$ means
\[
\limsup_{n\to\infty}
\sup_{\theta\leq\theta_0}E_\theta\varphi_n
\leq\alpha.
\]
This is stronger than merely checking
$E_{\theta_0}\varphi_n\to\alpha$.

Consistency against a fixed alternative $\theta>\theta_0$ means
\[
E_\theta\varphi_n\to1.
\]

\subsection{Wald and score statistics}

Let
\[
\sqrt n(\widehat\theta_n-\theta)
\xrightarrow{d}N(0,I(\theta)^{-1}).
\]
Two common Wald standardizations are
\[
W_n^{(0)}
=
\sqrt{nI(\theta_0)}(\widehat\theta_n-\theta_0),
\]
and
\[
W_n=
\sqrt{nI(\widehat\theta_n)}
(\widehat\theta_n-\theta_0).
\]
Under the null boundary and continuity of $I$, both converge to
$N(0,1)$.

The standardized score statistic is
\[
S_n=
\frac{\sum_{i=1}^ns_{\theta_0}(X_i)}
{\sqrt{nI(\theta_0)}}.
\]
It also converges to $N(0,1)$ under $\theta_0$.

A one-sided test rejects for a statistic exceeding
$z_{1-\alpha}$, where $z_u=\Phi^{-1}(u)$.

\subsection{Likelihood ratio statistics}

For testing the point null $\theta=\theta_0$ against a two-sided
alternative, define
\[
\Lambda_n=
2\{\ell_n(\widehat\theta_n)-\ell_n(\theta_0)\}.
\]
Under regularity and an interior null parameter,
\[
\Lambda_n\xrightarrow{d}\chi_1^2.
\]
Moreover,
\[
\Lambda_n=W_n^2+o_P(1)=S_n^2+o_P(1).
\]

For directional testing, use the signed likelihood root
\[
R_n=
\operatorname{sign}(\widehat\theta_n-\theta_0)
\sqrt{\Lambda_n}.
\]
Under suitable regularity,
\[
R_n=W_n+o_P(1)=S_n+o_P(1).
\]
A one-sided test rejects for $R_n>z_{1-\alpha}$.

The nonnegative likelihood-ratio statistic for a genuinely
inequality-constrained null has a boundary limit and should not be
confused with the ordinary two-sided $\chi_1^2$ statistic.

\subsection{Local alternatives and asymptotic power}

Fixed alternatives often give limiting power one, which does not
distinguish good tests. Consider instead
\[
\theta_n=\theta_0+\frac{h}{\sqrt n}.
\]
Under local uniform regularity,
\[
W_n\xrightarrow{d}
N\bigl(h\sqrt{I(\theta_0)},1\bigr)
\quad\text{under }P_{\theta_n}.
\]
Hence the limiting one-sided power is
\[
1-\Phi\bigl(z_{1-\alpha}-h\sqrt{I(\theta_0)}\bigr).
\]
Pointwise asymptotic normality at every fixed $\theta$ is not,
by itself, enough to justify this triangular-array statement.

The local likelihood expansion has the form
\[
\ell_n\left(\theta_0+\frac h{\sqrt n}\right)
-\ell_n(\theta_0)
=
h\Delta_n-\frac12h^2I(\theta_0)+o_P(1),
\]
where
\[
\Delta_n=\frac1{\sqrt n}\sum_{i=1}^ns_{\theta_0}(X_i)
\xrightarrow{d}N(0,I(\theta_0)).
\]
This local normal approximation explains both the common
first-order behavior of Wald, score, and likelihood-ratio tests
and the corresponding local power bound in regular models.

\section{Asymptotic Confidence Sets}

\subsection{Coverage and Wald intervals}

An asymptotic $(1-\alpha)$ confidence interval $C_n$ satisfies
\[
P_\theta(\theta\in C_n)\to1-\alpha
\]
for every fixed $\theta$. Uniform asymptotic coverage is a stronger
requirement.

If
\[
\frac{\sqrt n(\widehat\theta_n-\theta)}
{\sqrt{\widehat V_n}}
\xrightarrow{d}N(0,1),
\]
then a Wald interval is
\[
\left[
\widehat\theta_n-z_{1-\alpha/2}
\sqrt{\frac{\widehat V_n}{n}},
\quad
\widehat\theta_n+z_{1-\alpha/2}
\sqrt{\frac{\widehat V_n}{n}}
\right].
\]
For a regular efficient MLE,
$\widehat V_n=I(\widehat\theta_n)^{-1}$ is a standard choice.

\subsection{Inversion of tests}

If $\varphi_{n,\theta}$ tests the point null indexed by $\theta$,
its acceptance set defines
\[
C_n=\{\theta:\varphi_{n,\theta}=0\}.
\]
Thus
\[
P_\theta(\theta\in C_n)
=
1-P_\theta(\varphi_{n,\theta}=1).
\]
Score and likelihood-ratio confidence sets arise by this inversion.

Different constructions may be first-order equivalent while
behaving differently in finite samples. Inversion need not produce
an interval without additional shape or monotonicity conditions.
Boundary parameters and small samples can make naive Wald
intervals particularly unreliable.

\section{Posterior Asymptotics}

\subsection{Posterior distributions and two probability viewpoints}

Given a prior $\Pi$ and likelihood $L_n$, the posterior is
\[
\Pi(B\mid X_1,\ldots,X_n)
=
\frac{\int_B L_n(\theta)\,\Pi(d\theta)}
{\int_\Theta L_n(\theta)\,\Pi(d\theta)},
\]
when the denominator is positive and finite.

Conditional on the data, this is a probability distribution over
$\theta$. From a frequentist viewpoint, $\theta_0$ is fixed and the
posterior is a random probability measure because it depends on
the sample.

\subsection{Posterior consistency}

The posterior is consistent at $\theta_0$ if, for every
$\varepsilon>0$,
\[
\Pi\bigl(\|\theta-\theta_0\|>\varepsilon\mid X_1,\ldots,X_n\bigr)
\xrightarrow{P_{\theta_0}}0.
\]
Almost-sure posterior consistency replaces convergence in
probability by almost-sure convergence.

For a finite identifiable parameter space, positive prior mass at
every parameter value gives posterior consistency. Likelihood
ratios favor the true distribution, while the positive prior
weights do not decay with $n$.

In a continuous parameter space, $\Pi(\{\theta_0\})=0$ is normal
and does not prevent consistency. What matters is sufficient prior
support in neighborhoods of the truth, together with appropriate
model and separation conditions.

\subsection{Bernstein--von Mises principle}

In a regular finite-dimensional model, the posterior is
asymptotically normal around an efficient estimator. In one
dimension, the characteristic conclusion is
\[
\Pi\bigl(\sqrt n(\theta-\widehat\theta_n)\in\cdot
\mid X_1,\ldots,X_n\bigr)
\approx N(0,I(\theta_0)^{-1}).
\]

One sufficient abstract formulation uses:
\begin{enumerate}
\item Local quadratic likelihood expansion around
$\widehat\theta_n$, uniformly for bounded $t$:
\[
\ell_n\left(\widehat\theta_n+\frac t{\sqrt n}\right)
-\ell_n(\widehat\theta_n)
=
-\frac12I(\theta_0)t^2+o_P(1).
\]

\item A prior density continuous and positive at $\theta_0$,
and $\widehat\theta_n\xrightarrow{P}\theta_0$.

\item Tightness of the rescaled posterior:
for every $\varepsilon>0$,
\[
\lim_{M\to\infty}\limsup_{n\to\infty}
P_{\theta_0}\left(
\Pi\bigl(\sqrt n|\theta-\widehat\theta_n|>M\mid X\bigr)
>\varepsilon
\right)=0.
\]
\end{enumerate}
These conditions yield convergence to the normal law in total
variation, in $P_{\theta_0}$-probability.

The local expansion alone does not control posterior mass far
from the MLE; posterior tail control is an essential separate
ingredient. Ordinary posterior consistency alone does not
automatically supply the required $\sqrt n$-scale tightness.

\subsection{Sampling normality versus posterior normality}

The statements
\[
\sqrt n(\widehat\theta_n-\theta_0)
\xrightarrow{d}N(0,I(\theta_0)^{-1})
\]
and
\[
\Pi\bigl(\sqrt n(\theta-\widehat\theta_n)\in\cdot\mid X\bigr)
\approx N(0,I(\theta_0)^{-1})
\]
describe different objects.

The first concerns the distribution of an estimator under repeated
sampling with $\theta_0$ fixed. The second concerns a conditional
distribution over the parameter after the data have been observed.
Their matching first-order scales explain the agreement between
many Bayesian and frequentist procedures.

\subsection{Bayes estimators}

Let $\widetilde\theta_n$ be a posterior median. Posterior normality
and uniqueness of the limiting median imply
\[
\sqrt n(\widetilde\theta_n-\widehat\theta_n)
\xrightarrow{P}0.
\]
Thus a posterior median has the same first-order sampling
distribution as an efficient MLE.

A corresponding conclusion for the posterior mean requires
additional moment control. Weak or total-variation convergence of
posterior distributions alone does not imply convergence of their
means. A suitable condition is uniform integrability of the
rescaled posterior first moments.

Similarly, asymptotic statements about Bayes risks require control
of the relevant posterior loss moments, not merely convergence in
distribution.

\subsection{Credible intervals and confidence intervals}

A $(1-\alpha)$ credible set satisfies
\[
\Pi(C_n\mid X)=1-\alpha.
\]
This is a conditional posterior probability statement, not a
frequentist coverage statement.

Under Bernstein--von Mises conditions, equal-tailed posterior
intervals have endpoints
\[
\widehat\theta_n
\pm
\frac{z_{1-\alpha/2}}{\sqrt{nI(\theta_0)}}
+o_P(n^{-1/2}).
\]
If the centering estimator has the regular efficient sampling
distribution, these intervals have asymptotic frequentist coverage
$1-\alpha$.

First-order agreement does not imply exact finite-sample agreement.
Highest-posterior-density sets require additional density-level
regularity for analogous set-shape conclusions.

\subsection{Point-null testing}

A frequentist $p$-value and a posterior probability of a null
hypothesis are different quantities.

Under a purely continuous prior,
\[
\Pi(\{\theta_0\}\mid X)=0.
\]
A Bayesian test assigning positive probability to the point null
therefore requires an atom at $\theta_0$, for example a
spike-and-slab prior.

Even with such a prior, a small $p$-value need not imply a small
posterior null probability. This distinction is central to
Lindley-type phenomena.

\section{Logical Dependencies and Main Cautions}

The central logical chain is
\[
\begin{gathered}
\text{population identification}
+\text{uniform approximation}
\Longrightarrow\text{consistency},\\
\text{consistency}
+\text{local smoothness}
+\text{CLT}
\Longrightarrow\text{asymptotic linearity and normality},\\
\text{normality}
+\text{consistent scale estimation}
\Longrightarrow\text{asymptotic tests and intervals},\\
\text{additional regularity}
\Longrightarrow\text{information bounds and local optimality}.
\end{gathered}
\]

The following implications are generally false without additional
assumptions:
\begin{enumerate}
\item Identifiability of the model implies injectivity of a chosen
moment map.
\item A solution of the score equation is automatically a global MLE.
\item A unique population minimizer automatically gives consistency.
\item Convergence in distribution implies convergence of moments.
\item Pointwise asymptotic normality implies uniform validity or
validity under local alternatives.
\item The MLE is asymptotically optimal among all measurable estimators.
\item Posterior consistency alone implies Bernstein--von Mises.
\item Posterior normality alone implies convergence of posterior means
or risks.
\item A credible probability is the same as a confidence coverage
probability.
\end{enumerate}